\documentclass[../main]{subfiles}
\begin{document}
\chapter{Motor el\'ectrico{} sin escobillas}
Un motor el\'ectrico{} sin escobillas o motor \emph{brushless} es un motor el\'ectrico{} que no emplea escobillas para realizar el cambio de polaridad en el rotor.

Los motores el\'ectricos{} sol\'ian{} tener un colector de delgas o un par de anillos rozantes. Estos sistemas, que producen rozamiento, disminuyen el rendimiento, desprenden calor y ruido, requieren mucho mantenimiento y pueden producir part\'iculas{} de carb\'on{} que manchan el motor de un polvo que, adem\'as{}, puede ser conductor.

\img{images/brushless/brushless}{Coolers de ordenador desmontados.}

Los primeros motores sin escobillas fueron los motores de corriente alterna as\'incronos{}. Hoy en d\'ia{}, gracias a la electr\'onica{}, se muestran muy ventajosos, ya que son m\'as{} baratos de fabricar, pesan menos y requieren menos mantenimiento, pero su control es mucho m\'as{} complejo. Esta complejidad pr\'acticamente{} se ha eliminado con los controladores electr\'onicos{} de velocidad ESC (Por sus siglas en ingl\'es{} \emph{Electronic Speed Control}).

Otros motores sin escobillas, que s\'olo{} funcionan con corriente continua son los que se usan en peque\~nos{} aparatos el\'ectricos{} de baja potencia, como lectores de CD-ROM, ventiladores de ordenador, casetes, etc. Su mecanismo se basa en sustituir la conmutaci\'on{} (cambio de polaridad) mec\'anica{} por otra electr\'onica{} sin contacto. En este caso, la espira s\'olo{} es impulsada cuando el polo es el correcto, y cuando no lo es, el sistema electr\'onico{} corta el suministro de corriente. Para detectar la posici\'on{} de la espira del rotor se utiliza la detecci\'on{} de un campo magn\'etico{}. Este sistema electr\'onico{}, adem\'as{}, puede informar de la velocidad de giro, o si est\'a{} parado, e incluso cortar la corriente si se detiene para que no se queme. Tienen la desventaja de que no giran al rev\'es{} al cambiarles la polaridad (+ y -). Para hacer el cambio se deber\'ian{} cruzar dos conductores del sistema electr\'onico{}.

\subsection{Partes}
El motor cuenta con tres enrollados de cable de cobre con conexi\'on{} estrella, y dependiendo de cada tipo de motor, determinado n\'umero{} de polos para s\'i{} mismas. Com\'unmente{} los polos de las bobinas se encuentran en el estator y a su alrededor un n\'umero{} acorde a las bobinas de peque\~nos{} imanes que al energizarse cada una de las bobinas, o mejor dicho, cada par de bobinas hace girar al rotor para generar movimiento mec\'anico{}.

\subsection{Funcionamiento}
La manera de operar de estos motores llega a ser un poco diferente a cualquier otro motor de DC y AC. El principio del funcionamiento de los motores \emph{brushless} consiste en que, al momento de energizar dos polos de las tres bobinas que contiene, induzca un campo magn\'etico{} en las mismas para as\'i{} ser repelido por los imanes en el interior.

Al girar el rotor un paso hacia una direcci\'on{}, es al mismo tiempo repelido por un im\'an{} y atra\'ido{} por otro. Es ah\'i{} cuando se induce Potencial el\'ectrico{} en otra terminal del embobinado para soltarse de una de las previamente conectadas. Es ligeramente complicado controlar la velocidad de giro de este tipo de motor ya que es imposible hacer los cambios de conexiones entre las terminales de los embobinados a mano, es por ello que se utiliza un ESC (Controlador Electr\'onico{} de Velocidad, por sus siglas en ingl\'es{}) para poder variar las velocidades de giro por medio de Modulaci\'on{} por ancho de pulsos ya sea suministrado por un microcontrolador o por un transmisor de control remoto.

\subsection{ESC (Electronic Speed Controller)}
Es el dispositivo electr\'onico{} capaz de controlar el motor \emph{brushless} permitiendo hacer el intercambio de conexiones o polaridades en los embobinados, se controla PWM o pulsos.

El ESC cuenta con una serie de componentes electr\'onicos{} que en conjunto son capaces de hacer los movimientos necesarios para hacer funcionar el motor. Dependiendo del motor y de la potencia del mismo se pueden utilizar distintos tipos de controladores. Los m\'as{} comunes y comerciales son los de Circuitos integrados que puede hacerlo funcionar con una potencia baja pero con una facilidad de instalaci\'on{} y funcionamiento.

B\'asicamente{} est\'an{} compuestos por Transistores encargados de la etapa de potencia y a su vez sirven como interruptores de dos v\'ias{} que permiten alternar la polaridad. La otra parte importante es la de los \emph{Half-Bridge Drivers} que son los que reciben las respuestas de los pulsos para hacer el cambio de polaridad en los transistores. La parte m\'as{} importante en este tipo de controladores es la del controlador, ya que sin esta parte ser\'ia{} imposible hacer el movimiento del motor, no se puede hacer solamente mec\'anicamente{}, necesita una manipulaci\'on{} electr\'onica{}.

\subsection{Ventajas, desventajas y aplicaciones}
Los motores tienen la caracter\'istica{} de que no emplean escobillas en la conmutaci\'on{} para la transferencia de energ\'ia{}; en este caso, la conmutaci\'on{} se realiza electr\'onicamente{}. Esta propiedad elimina el gran problema que poseen los motores el\'ectricos{} convencionales con escobillas, los cuales producen rozamiento, disminuyen el rendimiento, desprenden calor, son ruidosos y requieren una sustituci\'on{} peri\'odica{} y, por lo tanto, un mayor mantenimiento.

Los motores \emph{Brushless} tienen muchas ventajas frente a los motores DC con escobillas y frente a los motores de inducci\'on{}. Algunas de estas son:
\begin{itemize}
\item Mayor relaci\'on{} velocidad-par motor.
\item Mayor respuesta din\'amica{}.
\item Mayor eficiencia.
\item Mayor vida \'util{}.
\item Menos ruido.
\item Mayor rango de velocidad.
\end{itemize}

\img{images/brushless/dron}{Dron impulsado con motores brushless}

Adem\'as{}, la relaci\'on{} par motor-tama\~no{} es mucho mayor, lo que implica que se pueden emplear en aplicaciones donde se trabaje con un espacio reducido.

El principal impedimento en la implementaci\'on{} de este tipo de motor es el coste del mismo ya que tiende a ser m\'as{} elevado que cualquier otro motor; el otro es el control del mismo, ya que como se mencion\'o{}, es imposible controlarlo manualmente por lo que necesita la ayuda electr\'onica{} para funcionar.

Actualmente, los motores \emph{Brushless} se emplean en sectores industriales tales como: Autom\'ovil{}, Aeroespacial, Consumo, M\'edico{}, Equipos de automatizaci\'on{} e instrumentaci\'on{}.

\subsection{Nomenclatura}
Cada motor tiene una numeraci\'on{} correspondiente que abarca sus caracter\'isticas{} esenciales, y que define de manera clara su comportamiento y capacidades. Ej: XXXX-XXXXKV.

Las primeras cuatro cifras se leen de dos en dos y corresponden a la longitud del motor y a su altura, respectivamente. Las cuatro siguientes cifras hacen menci\'on{} al coeficiente KV, que es una medida de la caracter\'istica{} f\'isica{} que define la calidad del motor.

El valor expresado en KV se refiere a la constante de revoluciones del motor, en resumen el n\'umero{} de RPM (Revoluci\'on{} por minuto) que ser\'a{} capaz de ofrecer el motor cuando se le aplique \qty{1}{V} (un voltio) de tensi\'on{}. No se debe confundir los KV del motor con kV (Prefijo ``kilo'' que expresa la multiplicaci\'on{} por 1000 por la unidad).

En los motores con KV bajos indica que el n\'umero{} de espiras es mayor, por lo tanto el hilo de cobre es m\'as{} fino. El total de amperios que circulan por el motor es inferior a otros con KV m\'as{} alto. Se utilizan en drones que necesitan mucho par y poca velocidad.

En los motores con KV altos indica que el n\'umero{} de espiras es menor, por lo tanto el hilo de cobre es m\'as{} grueso. El total de amperios que circulan por el motor es superior a otros con KV m\'as{} bajo. Se utilizan en drones de carreras, aparatos que necesitan poco par y mucha velocidad.

\video{https://www.youtube.com/watch?v=NnUiAgUundw}{Funcionamiento del motor sin escobillas}

\actividad{Leer, contestar e investigar.}
\begin{enumerate}
\item \textquestiondown{}Qu\'e{} es un motor el\'ectrico{} sin escobillas?
\item \textquestiondown{}Cu\'ales{} son las desventajas de los motores el\'ectricos{} con escobillas?
\item \textquestiondown{}Cu\'al{} es la funci\'on{} de un controlador electr\'onico{} de velocidad (ESC) en un motor \emph{brushless}?
\item \textquestiondown{}C\'omo{} se realiza la conmutaci\'on{} en un motor \emph{brushless}?
\item \textquestiondown{}Qu\'e{} ventajas ofrecen los motores \emph{brushless} en comparaci\'on{} con los motores de corriente continua y los motores de inducci\'on{}?
\item \textquestiondown{}Cu\'al{} es la ventaja principal de los motores \emph{brushless} en t\'erminos{} de relaci\'on{} par motor-tama\~no{}?
\item \textquestiondown{}Cu\'al{} es la funci\'on{} del coeficiente KV en la numeraci\'on{} de un motor \emph{brushless}?
\item \textquestiondown{}Qu\'e{} indica el valor de KV en un motor \emph{brushless}?
\item \textquestiondown{}Cu\'ales{} son las aplicaciones industriales de los motores \emph{brushless}?
\item \textquestiondown{}Cu\'al{} es la diferencia entre los motores \emph{brushless} con KV alto y KV bajo en t\'erminos{} de espiras y rendimiento?
\end{enumerate}
\end{document}